\scanpage{95}
% archive foliation: 94
\uline{Morphisme de type fini} \deleted{[\unsure{prés. finie|quasi-fini}]} d'ens schématiques : tel qu'il existe une partition de $X$ en ensembles loc. fermés constructibles affines $X_i$, telle que les $f\,|\,X_i$ soient des morphismes de type fini \deleted{[\unsure{prés. finie|quasi-fini}]}.

\margin{\unsure{cela implique} \unsure{la définition} \unsure{présentation finie} ?}

Un tel morphisme d'ens schématique est isom. à un morphisme de type fini de schémas affines, (considéré comme morphisme d'ens schématiques). \uline{Transitivité}.

Ens. schématique sur un autre, de type fini sur un autre. Tous sont sur $\unsure{\mathbf Z|Z}$ !

Ens. schématique défini par une partie constructible d'un ens. schématique, son immersion canonique, décomposition de $X$ en somme $X=E\amalg(X-E)$, critère de factorisabilité par une partie constructible; \uline{universalité} !

Produit fibré [et $\varprojlim$ finies] dans la catégorie des ens schématiques. \uline{Existe}, et le foncteur \deleted{préschémas}
\diagram{\begin{tikzcd}
(\text{présch qu compt qu sép})\arrow[r]&(\text{ens schém})
\end{tikzcd}}
commute aux produits fibrés.

Si on a $X$ et $Y$ sur $Z$, si $X',Y'$ sont des parties constructibles de $X,Y$ alors
\[
X'\times_ZY'=\operatorname{pr}_1^{-1}(X')\cap\operatorname{pr}_2^{-1}(Y')\quad\ldots.
\]
% Reversed blue typewriting is transcribed on its printed face, scan 96.
